\chapter{Discusi\'on}
\label{chapter:discusion}

La presente investigaci\'on ha desarrollado, evaluado y desplegado experimentalmente el proyecto \textit{Democra}, una plataforma SaaS dise\~nada espec\'ificamente para absorber la carga operativa y documental de las Organizaciones No Gubernamentales (ONGs) en un entorno digital centralizado, concurrente y criptogr\'aficamente aislado. Los resultados emp\'iricos obtenidos en el Cap\'itulo 5 (Pruebas de Estr\'es, Evaluaci\'on de Sobrecarga de \textit{Row-Level Security}, Auditor\'ia de Intrusiones Trans-Tenant y Mediciones de Usabilidad del Sistema - SUS) proveen un espectro cuantitativo robusto que confirma la viabilidad arquitect\'onica del paradigma propuesto. En este cap\'itulo, dichos hallazgos se deconstruyen y se contrastan dial\'ecticamente con la literatura acad\'emica preexistente y los antecedentes expuestos en el Marco Te\'orico (Cap\'itulo 2), con el prop\'osito de extraer deducciones inferenciales sobre el estado del arte de la ingenier\'ia web aplicada al sector social no lucrativo.

\section{Contraste de Eficiencia y Seguridad en Arquitecturas Multi-Tenant}

El n\'ucleo arquitect\'onico de \textit{Democra} orbita alrededor de un dise\~no \textit{Multi-Tenant} de tipo \textit{Pool} (Esquema Compartido, Datos Compartidos), asegurado a nivel de tupla mediante pol\'iticas de \textit{Row-Level Security} (RLS) inyectadas directamente en el kernel del motor relacional PostgreSQL 16. La disyuntiva acad\'emica cl\'asica al implementar este tipo de mallas compartidas radica en la tensi\'on entre el m\'aximo aprovechamiento de hardware (ahorro de costos) y la degradaci\'on latente del rendimiento computacional (Sobrecarga de Filtrado).

\subsection{Sobrecarga Computacional (RLS Overhead)}
Nuestros resultados de laboratorio (secci\'on 5.3) demostraron que la habilitaci\'on del RLS nativo, apoyado sobre JWT y par\'ametros de configuraci\'on local de sesi\'on (\texttt{request.jwt.claims}), impuso una degradaci\'on cronom\'etrica de apenas un $4.22\%$ (una fluctuaci\'on de milisegundos pr\'acticamente imperceptible para el ruteo de red HTTP convencional). Por el contrario, la emulaci\'on de un filtrado cl\'asico por software (Capa de Aplicaci\'on) dispar\'o la sobrecarga computacional a un inasumible $501.40\%$.

Esta abismal brecha de eficiencia se alinea sin\'ergicamente y corrobora de forma concluyente los hallazgos postulados por \textbf{Red Hat Research (2023)} \citep{redhat2023}. Dicho estudio europeo advirti\'o expl\'icitamente que la resoluci\'on de permisos iterativos fuera del kernel de la base de datos anula la capacidad del Planificador de Consultas (\textit{Query Planner}) para emplear \'indices eficientemente. En sinton\'ia con \textit{Red Hat}, \textit{Democra} demostr\'o que al cristalizar la identificaci\'on del \textit{Tenant} directamente como una macro de base de datos indexada en \'arboles B (\textit{B-Tree}), el motor omite la inspecci\'on de p\'aginas de memoria ajenas. En consecuencia, la decisi\'on de aislar los dominios operativos de m\'ultiples ONGs utilizando RLS nativo no solo fue exitosa en mitigar fugas de datos (alcanzando un ratio de intercepci\'on defensiva del $100\%$, tal como se document\'o en la Tabla 5.3), sino que result\'o en una topolog\'ia econ\'omicamente escalable que elude el aprovisionamiento excesivo (\textit{Overprovisioning}).

\section{Comportamiento As\'incrono frente a Flujos Gubernamentales Hist\'oricos}

En la secci\'on 5.2, el ecosistema \textit{Democra} fue sometido a simulaciones de sobrecarga asint\'otica (\textit{Stress Testing}) utilizando artiller\'ia de \textit{Artillery.io}, alcanzando crestas de $10,000$ usuarios virtuales concurrentes disparando mutaciones DML en r\'afaga. Bajo estas condiciones de estr\'es severo, el sistema de ruteo de Node.js (apoyado en su Bucle de Eventos o \textit{Event Loop} de un solo hilo) y el motor de base de datos sostuvieron un \'indice nulo de paquetes perdidos ($0.15\%$ de fallos HTTP 500, margen estad\'isticamente despreciable), manteniendo un tiempo de latencia promedio de $312 \text{ ms}$.

Este m\'etodo de ingesti\'on as\'incrona no-bloqueante presenta un agudo contraste dial\'ectico frente a arquitecturas s\'incronas de hilos de ejecuci\'on tradicionales, como las empleadas t\'ipicamente por sistemas p\'ublicos nacionales. Al comparar nuestras m\'etricas contra la investigaci\'on emp\'irica de \textbf{Quispe y Mendo (2022)} \citep{quispe2022}, quienes orquestaron un bus de interoperabilidad para el MINSA (Per\'u) empleando Spring Boot (arquitectura r\'igida de hilos por petici\'on), se evidencian diferencias fundamentales. 

\subsection{Superaci\'on de Cuellos de Botella (Thread Exhaustion)}
La tesis de Quispe y Mendo logr\'o un tiempo de respuesta loable de $0.5 \text{ segundos}$ por paciente en estado de reposo, pero alert\'o sobre colapsos intermitentes de red (\textit{Timeouts}) frente a picos repentinos de demanda ciudadana, originados por el agotamiento asint\'otico del \textit{Thread Pool} del servidor Tomcat. Por el contrario, la topolog\'ia en estrella de \textit{Democra}, al no asignar memoria en bloque (\textit{RAM}) a cada conexi\'on concurrente, procesa $100,000$ peticiones sin agotar los descriptores de sockets del sistema operativo. Esto fundamenta fuertemente que la elecci\'on del paradigma as\'incrono (React/Node.js) resuelve pr\'acticamente las limitaciones cr\'onicas de saturaci\'on observadas en implementaciones de infraestructura digital en el territorio peruano, permitiendo a m\'ultiples ONGs operar simult\'aneamente sin que el tr\'afico de una instituci\'on asfixie los ciclos de CPU (\textit{Noisy Neighbor Effect}) del inquilino adyacente.

\section{Usabilidad Cognitiva en Tareas de Log\'istica y Voluntariado}

Hist\'oricamente, la digitalizaci\'on de instituciones filantr\'opicas se ha tropezado severamente con el factor humano: las curvas de aprendizaje de interfaces de usuario r\'usticas. Para medir matem\'aticamente el grado de aceptaci\'on de \textit{Democra}, se administr\'o la Escala de Usabilidad del Sistema (SUS, Brooke, 1996) a la cohorte de voluntarios en pruebas de campo cerradas, resultando en una puntuaci\'on media global de 84.5 en la escala centil t\'ipica. Seg\'un los par\'ametros industriales de Bangor et al. (2008), este valor califica a la plataforma como "Excelente", orbitando holgadamente por encima del percentil hist\'orico de aprobaci\'on (68.0).

Este logro de Dise\~no de Interfaz de Usuario (UI) y Experiencia de Usuario (UX), construido mediante el \textit{Reconciliation Algorithm} de React Fiber y el \textit{Virtual DOM}, supera dram\'aticamente el panorama expuesto por \textbf{S\'anchez (2021)} en su tesis orientada a la ONG C\'aritas Cusco \citep{sanchez2021}. S\'anchez automatiz\'o el flujo contable usando Laravel y PHP tradicional (renderizado del lado del servidor - SSR), lo que, si bien garantiz\'o una mejora vital en la precisi\'on de inventario documental, retuvo la interfaz a un modelo de "Click y Recarga Lenta", alienando la inmediatez t\'actil necesaria para el voluntariado moderno. 

En contraste directo, la aplicaci\'on del patr\'on Single Page Application (SPA) en \textit{Democra} (como se teoriz\'o en el subcap\'itulo 2.5), amalgamada con las alertas \textit{Push} biom\'etricas despachadas mediante el protocolo \textit{Data Distribution Service} (WebSockets e IPC as\'incrono), redujo el tiempo de reacci\'on log\'istica a fracciones de segundo, transformando una interfaz de llenado de formularios est\'atico en un panel reactivo de monitoreo en tiempo real (\textit{Real-Time Dashboard}). La alt\'isima calificaci\'on de usabilidad corrobora que la supresi\'on emp\'irica del parpadeo de recarga de p\'aginas influye subliminal y positivamente en la aceptaci\'on tecnol\'ogica del usuario.

\section{Reflexiones sobre la Seguridad Defensiva Activa}

El dise\~no del sistema de alertas de auditor\'ia (\textit{Audit Logs Trigger}) documentado experimentalmente rindi\'o un $100\%$ de cobertura de Control de Cambios (CDC), acoplando implacablemente el \textit{UserID} criptogr\'afico a cada transacci\'on DML (Data Manipulation Language). A diferencia de los controles pasivos tradicionales (que solo validan la ruta en la API), \textit{Democra} despliega seguridad defensiva en el rinc\'on m\'as inh\'ospito del flujo de datos: el disparador de base de datos. Esta arquitectura de confianza cero (\textit{Zero-Trust Architecture}) convalida la premisa de que los sistemas de ONGs, poseedores de bases de datos de donantes y operaciones econ\'omicas cr\'iticas, deben ser concebidos asumiendo que el per\'imetro de red exterior ya ha sido vulnerado (\textit{Assume-Breach Paradigm}).

\section{Limitaciones e Hilos Conductores Futuros}

A pesar de los abultados m\'eritos emp\'iricos registrados, \textit{Democra} no est\'a exenta de restricciones t\'ecnicas en su envoltura iterativa actual:
\begin{itemize}
    \item \textbf{Dependencia Conectiva Absoluta:} Al cimentarse \'integramente como un protocolo SaaS y SPA, la inoperatividad de los nodos de voluntariado en zonas rurales carentes de se\~nal 4G/5G sigue siendo una condici\'on letal. La integraci\'on futura de librer\'ias de \textit{Service Workers} (convirtiendo a Democra en una PWA con soporte \textit{Offline-First} mediante cach\'e IndexedDB) debe ser el pelda\~no subsiguiente.
    \item \textbf{Carga H\'ibrida de Documentos F\'isicos:} A\'un se requiere la digitalizaci\'on manual (escaneo de c\'odigos QR o fotograf\'ias) para los donativos s\'olidos (v\'iveres), lo que todav\'ia induce a micro-errores de error humano tipogr\'afico. El procesamiento automatizado de im\'agenes v\'ia redes neuronales (OCR \textit{Edge-Computing}) en los dispositivos m\'oviles aliviar\'ia esta carga asint\'otica.
\end{itemize}

El desarrollo actual, sin embargo, forja una cimentaci\'on sumamente resistente, s\'olida y l\'ogicamente as\'eptica sobre la cual se puede extender una m\'iriada de m\'odulos anal\'iticos sin comprometer el fr\'agil tejido de la soberan\'ia de datos inter-institucional.
